\chapter{From Insulin to Gene Therapy}

\begin{kaichang}
A pharmacy refrigerator holds small white boxes containing something like a marker pen: an insulin pen. Hundreds of millions of people worldwide depend on it, inserting a needle into their abdomen daily and pressing a button. A few microliters of clear liquid enter the body, bringing blood sugar under control.

Guess three things. First, where does that insulin come from? A century ago the answer might have been ``extracted from two hundred pigs' pancreases''; today's answer may surprise you. Second, if illness results from a wrong letter in the body's blueprint, could doctors correct the blueprint instead of requiring daily injections? Third, is correcting it the same as correcting it before it passes to children?

Write down your guesses. By the chapter's end, these questions will connect modern medicine's most exciting path with its greatest need for clear thinking.
\end{kaichang}

\textbf{Translator safety note:} The opening injection description is not dosing instruction. Insulin type, concentration, dose, and delivery must follow the individual's prescribed plan; do not convert the example's liquid volume into a dose or change treatment yourself. See \href{https://www.niddk.nih.gov/health-information/diabetes/overview/insulin-medicines-treatments}{NIDDK insulin and diabetes-treatment guidance}.

\section{The Protein Inside a Pen}

Start with what is observable. In type 1 diabetes, the immune system damages insulin-producing pancreatic-islet cells, interrupting supply. Insulin is an instruction to open doors: when blood glucose rises, it tells muscle, liver, and fat cells to open membrane glucose channels and store sugar. Without it, sugar accumulates in blood while cells starve, producing thirst and weight loss; prolonged high glucose damages vessels, kidneys, and nerves.

Type 2 diabetes is more common and differs: insulin is still produced, but cells ignore the instruction, called insulin resistance. Detailed medicine lies beyond this chapter. Remember that the missing item is not sugar but a molecular key.

The key is a substance: insulin is a protein of 51 amino acids in two disulfide-linked chains, folded into a compact shape. Chapter 60 described amino-acid chains becoming molecular machines. As a substance, it faces our five familiar questions: what comprises it? Amino acids. How are they connected? As specified by the gene's base sequence. Can it change, and how quickly? Yes, if we find a machine that reads its blueprint. Life's best blueprint-reading machines are cells.

The next story is therefore not how chemists synthesize insulin but how we recruit other cells to produce it.

\section{From Pancreases to Bacterial Factories}

When insulin was discovered in 1922, slaughterhouses were its only source. Grinding pig and cattle pancreases, extracting with acid, and purifying produced small quantities of crystals. Animal insulin differs from human insulin by one or two amino acids; most patients could use it, but some developed allergies, and production never met demand.

The turning point came when Chapter 82's tools matured: restriction enzymes cut DNA, ligase joins it, and plasmids carry genes into bacteria. Insert a human insulin gene into a plasmid and E. coli, and each division copies the blueprint while the cells produce human insulin. A fermentation tank becomes a miniature pharmaceutical factory.

\begin{zhengju}
The real story is more interesting than the legend. In 1978, two-year-old Genentech and City of Hope laboratories in Los Angeles announced E. coli production of human insulin. They \emph{did not} isolate the insulin gene directly from human cells, then technically difficult. Instead, knowing both amino-acid sequences, they worked backward to suitable DNA sequences, chemically assembled the genes piece by piece, expressed them in separate bacterial populations, and joined the two chains outside the cells.

Why the detour? Deriving a blueprint from a protein tests the central dogma's credibility (Chapter 64): if base sequence determines amino-acid sequence, a gene written backward from the code table should produce the correct protein in bacteria. It did, closing the chain. In 1982, recombinant human insulin became the first approved genetically engineered medicine. Challenges remained: manufacturers had to show identity with natural insulin and absence of bacterial fever-causing contaminants, using successive chromatography and rigorous testing. Today, engineered yeast has replaced E. coli as the main producer, yielding higher purity.
\end{zhengju}

Producer bacteria or yeast grow in giant stainless-steel \textbf{bioreactors}. Unlike high-temperature, high-pressure chemical vessels, these cater to living cells: roughly $37\,^{\circ}\mathrm{C}$ for bacteria, lower for yeast; automated acid--base pH control; precisely formulated nutrients down to trace elements; and oxygen supplied by mixing and aeration. Mistreat these living workers and they stop, die, or degrade the product. Yields reach grams per liter, so a few hundred liters can replace a slaughterhouse's month of pancreases.

\begin{qiaoliang}
Consider the scale change. In the 1960s, extracting one kilogram of insulin crystals required processing over ten tons of pancreases from more than a hundred thousand pigs. Modern yeast fermentation can yield several grams per liter. A 1,000-liter vessel at three grams per liter produces about three kilograms per batch. A small fermentation facility's output over days can equal hundreds of thousands of pigs. Animal insulin rapidly withdrew after engineered products arrived because the improvement was not merely quality but orders of magnitude in quantity. Such scale changes can transform society.
\end{qiaoliang}

\section{The Protein-Medicine Family}

Insulin opened a door to a large family. If microbes can make proteins from blueprints, in principle any useful protein can be produced once its gene is known.

Growth hormone formerly came from deceased donors' pituitaries, with dozens of bodies supplying one child's annual needs and occasional transmission of fatal prion disease. Recombinant products ended both nightmares. Clotting factor VIII for hemophilia formerly came from pooled donor plasma, sometimes contaminated with hepatitis viruses and HIV during the 1980s. Recombinant versions transformed treatment.

The family's most intricate members are \textbf{monoclonal antibodies}. Recall Chapter 70's Y-shaped immune proteins with arms binding particular antigens. Natural immunity produces mixtures of thousands of antibodies. Fuse a B lymphocyte producing one antibody with a myeloma cell, however, and the hybridoma proliferates indefinitely while producing just that antibody. This 1975 invention later earned a Nobel Prize. Modern engineering directly introduces antibody genes into production cells, even bypassing mice. Monoclonal antibodies target tumor markers, block inflammatory signals, and neutralize viruses, forming one of medicine's fastest-growing categories.

Vaccines are becoming molecular too. Traditional platforms culture viruses before inactivation or attenuation, taking years. New platforms instead provide not virus but its component blueprint: mRNA encoding a viral surface protein, delivered in lipid nanoparticles. The body's cells temporarily make viral protein for immune rehearsal, following Chapter 64's translation chain. The mRNA is then degraded without entering the nucleus or rewriting the genome. During COVID-19 in 2020, this reduced vaccine design to weeks after receiving the viral sequence.

\begin{shiyan}
Try the logic of a bioreactor at home.

\textbf{Materials}: dry yeast, sugar, three clean mineral-water bottles, three balloons, warm water, and a marker.

\textbf{Procedure}: fill each bottle halfway with warm water, add a spoonful of sugar, and mix. Bottle A should be about $35\,^{\circ}\mathrm{C}$, warm but not hot to the touch; B uses cold tap water; C uses hotter water, about $55\,^{\circ}\mathrm{C}$, handled by a parent to prevent scalds. Add half a spoonful of yeast to each, mix, and fit balloons over the openings. Leave them in place and record balloon inflation every ten minutes for at least forty minutes.

\textbf{What to observe}: which balloon inflates fastest? Yeast enzymes convert sugar into carbon dioxide and ethanol through Chapter 58's fermentation. A larger balloon indicates a busier cell factory. One temperature range is just right: cold slows enzymes; excessive heat denatures proteins and kills workers. Hence bioreactors must control temperature and pH precisely.

\textbf{Safety reminder}: all materials are food grade; the only risk is hot water, which parents should handle. Pour the liquid down the drain afterward and do not drink it.
\end{shiyan}

\textbf{Translator safety note:} Hot water is not the only hazard: uninflated balloons and fragments can choke children, and latex can trigger allergy. Adults should handle hot water using a thermometer, not a hand-temperature test, and keep balloons out of mouths and away from young children. Do not replace the balloon with a rigid sealed cap. See \href{https://www.cpsc.gov/s3fs-public/5098.pdf}{CPSC scald precautions}, \href{https://www.cpsc.gov/s3fs-public/5087.pdf}{CPSC balloon warning}, and \href{https://www.cdc.gov/niosh/docs/2012-119/default.html}{NIOSH latex guidance}.

\section{Repair the Blueprint, Not the Part}

All these medicines \textbf{supplement} missing proteins rather than \textbf{repairing} the underlying defect. Patients need lifelong injections. Could we instead modify the faulty gene directly?

Chapter 82 explained CRISPR targeting, cutting, and rewriting. With the tools available, engineering questions remain: which cells receive them, how much, and with what outcome? Here lies one of the book's most important boundaries:

\begin{itemize}[itemsep=2pt]
\item \textbf{Somatic gene therapy} changes genes in some of a patient's ordinary cells, such as retinal cells for inherited blindness or liver cells for a metabolic disease. Changes end with the person and \textbf{do not pass to children}. Conceptually, this resembles replacing some parts; several therapies are approved worldwide.
\item \textbf{Germline editing} changes sperm, eggs, or early embryos. The change enters \emph{every cell} of the future child, including reproductive cells, and \textbf{passes down generations}.
\end{itemize}

The difference concerns responsibility as well as reach. A somatic-treatment patient accepts risks through informed choice. Germline consequences reach people not yet alive and unable to consent. An off-target error incorporated there resembles a typo in a printing master, reappearing in every later copy. In 2018, a researcher announced the birth of two babies after embryo editing intended to confer HIV resistance. The gene's full functions, unintended editing damage, and long-term health outcomes were uncertain. Scientists almost unanimously condemned the act, and the researcher was imprisoned. The prevailing consensus is that germline editing is not ready for clinical use. This boundary is scientific honesty, not hostility to science: unknown consequences do not justify irreversible intervention.

\section{Turning Immune Cells into Hunters}

One gene therapy scarcely needs viral vectors: it removes cells, modifies them, and returns them. This is CAR-T therapy.

Recall Chapter 70: T cells are immune hunters whose surface receptors recognize suspicious molecules and attack cancer cells. Tumors can disguise themselves. CAR-T takes a patient's T cells and genetically equips them with a chimeric antigen receptor, or CAR, designed for a marker on that cancer. The modified cells expand to hundreds of millions in culture before infusion. It resembles taking police out for special training, fitting recognition glasses, and returning them to the city.

For some leukemias and lymphomas, the results have been remarkable: patients with exhausted treatment options achieved prolonged remission, and some of the earliest pediatric recipients have remained disease-free for over ten years. But remarkable outcomes require a full accounting:

\begin{itemize}[itemsep=2pt]
\item \textbf{Efficacy has limits}. The strongest CAR-T results concern several blood cancers. Solid tumors such as lung and liver cancers lack sufficiently specific markers and have immune-suppressive environments, limiting success.
\item \textbf{Risks are real}. Activated cells, numbering hundreds of millions, can release signaling molecules at once, causing a cytokine storm with fever, low blood pressure, and potentially life-threatening illness. Neurological injury is also possible. Experienced hospitals must monitor recipients; this is no ordinary injection.
\item \textbf{Cost and access matter}. Individually manufactured cell batches take weeks and have cost millions of yuan. Most countries currently reserve treatment for patients whose other therapies failed, with differing insurance coverage.
\item \textbf{Time has not finished its work}. The first CAR-T product was approved only in 2017. How long do modified cells persist, and could problems appear decades later? Long follow-up alone can answer. Gene-therapy recipients therefore commonly enter registries lasting more than ten years. Approval begins observation rather than ending the story.
\end{itemize}

\textbf{Translator safety note:} CAR-T requires specialist assessment and product-specific monitoring; it is not suitable for every cancer or patient. In addition to acute immune and neurological toxicity, regulators warn about secondary T-cell malignancies and the need for long-term monitoring. Do not seek unregulated cell or gene treatment. See \href{https://www.fda.gov/vaccines-blood-biologics/safety-availability-biologics/fda-investigating-serious-risk-t-cell-malignancy-following-bcma-directed-or-cd19-directed-autologous}{FDA CAR-T safety information} and \href{https://www.fda.gov/vaccines-blood-biologics/cellular-gene-therapy-products/information-about-self-administration-gene-therapy}{FDA gene-therapy cautions}.

\begin{wuqu}
``Gene therapy is one injection that fixes all bad genes throughout the body permanently.'' Both ideas are mistaken. Delivery cannot evenly reach tens of trillions of cells; current therapies mostly access a fraction within particular organs such as eyes, liver, or blood, leaving edited and unedited cells together in a mosaic. Permanence also ignores uncertainty: editing may go off target, inserted genes may be silenced, and benefits may fade. Real gene therapy resembles an expensive, risky major repair followed by lifelong observation, not magic.
\end{wuqu}

\section{Precision Medicine Is Not Fortune-Telling}

One shared conviction connects this chapter: reading personal genomes (Chapter 81), editing them (Chapter 82), and producing proteins from instructions could replace one-drug-for-everyone medicine with molecularly tailored care. This is \textbf{precision medicine}.

Some promises have been realized. Certain lung cancers carry mutations targeted by drugs far more effective than traditional chemotherapy for those patients; testing before treatment selection is routine. Genetic differences in liver enzymes also alter drug metabolism and appropriate doses, the field of pharmacogenomics. Chapter 85 examines what testing can and cannot establish.

Remember two speed bumps.

First, \textbf{most common diseases are not single-gene disorders}. Hypertension, diabetes, and depression involve hundreds or thousands of small genetic effects combined with diet, activity, sleep, and luck. Predictions are probability distributions, not verdicts. A high-risk genotype may never lead to disease, while someone without it becomes ill. Like Chapter 76's population heritability, statistics describe group differences, not an individual's script.

Second, \textbf{evidence comes in levels}. One miraculous cure is weak evidence, potentially reflecting placebo effects, fluctuating disease, or survivor bias. Case series, controlled studies, and randomized trials follow, topped by syntheses of multiple randomized trials. Treatments spend years climbing from laboratories to guideline recommendations. On hearing breakthrough or cure, ask which evidence level has been reached. For individual decisions, that standard outweighs any book. This book explains principles, not a substitute for clinicians and proper care.

\begin{tiaozhan}
Why might even precise rare-disease treatments often fail? Suppose tools reach 10\% of target-organ cells and correctly repair 90\% of those, with 10\% failing through off-target effects or other causes. The repaired fraction is $0.1\times0.9=9\%$. If 20\% functional cells are needed for symptom improvement, treatment misses the threshold despite apparently favorable individual steps. Now consider germline risk: even an accidental-edit probability of one in ten thousand per site cannot be ignored when an embryo has billions of sites and tools scan many positions. The posterior probability of one accident is no longer negligible, and that accident passes to all descendants. Low probabilities multiplied by large numbers underpin this chapter's caution. You may skip this section, but it merits rereading.
\end{tiaozhan}

\section{When These Treatments Fail}

As always, even a good model needs clear boundaries. Reading and writing life has several.

First, it handles \textbf{single genes, known mechanisms, and reachable target organs} best: replace or repair one missing or defective protein. Coronary disease, Alzheimer's disease, and most cancers instead involve many genes, steps, and time-dependent system changes, not one faulty blueprint. Extending single-gene success to every illness is a common overreach.

Second, immunity helps and hinders. Unrecognized vectors and foreign proteins can trigger attacks, nullifying therapy or causing dangerous inflammation. In 1999, teenager Jesse Gelsinger died from a severe immune reaction in an early gene-therapy trial, prompting years of field-wide suspension and reassessment. The costly lesson was that delivery toxicity is a primary scientific question, not an engineering afterthought.

Third, production-line success does not guarantee bodily efficacy. Protein medicines are vulnerable to heat and stomach acid, so insulin still requires injection rather than oral use: digestion treats it as ordinary food, regardless of value (Chapter 63). However promising in a dish, a molecule must negotiate absorption, distribution, metabolism, and excretion before producing benefit.

Fourth, therapies inhabit society even when their science works. Prices, capacity, insurance, cold chains, and regional care determine how many people a paper breakthrough reaches. An extraordinary drug available to few may contribute less to overall health than an affordable reliable one. Molecular precision and access belong on the same scale when evaluating progress.

\section{Back to the Pen}

We can now answer the opening questions.

Where does pen insulin come from? A cell factory: human insulin genes reside in yeast or bacterial plasmids, while tanks regulate temperature, pH, nutrients, and oxygen. Behind a few microliters lies a DNA-to-protein production line. Did you guess correctly?

Can we repair blueprints rather than parts? Yes, for now in somatic cells of reachable organs under strict oversight. CAR-T and several gene therapies have saved people with formerly untreatable conditions. Germline editing before transmission to children crosses the mainstream scientific boundary; those questions have different answers.

The pen itself continues evolving: modified amino acids produce faster or longer-lasting insulin analogues by altering aggregation beneath the skin and absorption rates. Structure determines properties, Chapter 60's rule still operating in pharmacies. Combined with continuous glucose monitoring, this brings patients closer to an artificial pancreas. Yet dose, evidence, individual variation, and time always separate molecules from bedsides.

\section{Where the Chain Leads}

Here we saw reading and writing life in medicine. Applied elsewhere, the tools raise different debates: is crop engineering breeding or risk-taking? Can copying a genome copy a person? Where are the boundaries of designing cells from scratch? Chapter 84 examines genetic modification, cloning, and synthetic biology. Chapter 85 instead considers reading without changing genes, asking what a report can and cannot tell you. Sharper tools make whether to use them a question for everyone, not technology alone.

\section*{Try It Yourself}

\begin{sikaoti}
\item Draw the information flow from a human insulin gene to pharmacy insulin. Identify the information's form at each step, DNA, mRNA, or protein, and the machine performing it. State that this is a representation, with no tiny people reading blueprints in the tank.
\item Chapter 82 says enzymes lower activation energy without determining reaction direction. Apply this here: do bioreactor cells supply materials, energy, or instructions? What happens if any one is missing?
\item Why might pig-derived insulin cause allergy in some patients while recombinant human insulin rarely does? Explain through protein structure and immune recognition. (Hint: Chapter 70.)
\item Explain to an older relative, without technical terms if possible, why editing blood-cell genes to treat illness differs fundamentally from editing an embryo.
\item A report claims a new drug has a 90\% cure rate. List at least three questions, such as compared with what, how many patients, and how long the follow-up lasted, and explain the alternative explanation each addresses.
\item A somatic gene therapy delivers tools to 15\% of liver cells and correctly repairs 80\% of those. Clinicians estimate 10\% repaired cells would substantially improve symptoms. Using the challenge section's method, could the approach cross the threshold? If the result were close to the threshold, which number would you recommend improving first?
\end{sikaoti}
